\chapter{Workflow-Blöcke}\label{kap:workflow-bloecke}\index{Workflow-Block}

Dieser Anhang beschreibt alle Blocktypen, die der Node-Workflow (Kapitel~\ref{kap:workflow})
in der vorliegenden Version kennt: die eingebauten Typen der Workflow-Maschine (RTT-Fusion) und die
Typen, die die Plugins \plugin{Raytracer}, \plugin{FourierGrating} und \plugin{hp-FEM} liefern. Die Angaben stammen
aus der Blockbibliothek, die die Workflow-Maschine dem Editor übergibt; im Editor finden Sie
die Typen unter \gui{Add block}, geordnet nach Kategorie -- dort aber nur die, die Ihre
Installation rechnen kann (Abschnitt~\ref{sec:bloecke-auswertung}).

\section{Lesehilfe}\label{sec:bloecke-lesehilfe}

Zu jedem Block stehen:

\begin{description}
  \item[Kategorie] die Kategorie (bestimmt die Farbe der Kopfzeile und das Untermenü von
        \gui{Add block}): \code{source}, \code{domain}, \code{splitter}, \code{detector},
        \code{conversion}, \code{data}, \code{analysis}, \code{meta}.
  \item[Anbieter] das Plugin, das den Typ liefert und rechnet. \enquote{Workflow (RTT-Fusion,
        eingebaut)} heißt: Die Workflow-Maschine rechnet ihn selbst.
  \item[Beschreibung] die Kurzbeschreibung der Bibliothek (englisch, kursiv; sie erscheint auch
        als Tooltip in \gui{Add block}) und eine Erläuterung.
  \item[Licht und Daten] alle Anschlüsse (\emph{sockets}), getrennt nach Licht und Daten, je
        als \textit{Name} (\textit{Richtung}, \textit{Darstellungen} bzw.\ \textit{Datentyp},
        ggf.\ \textit{Seite}, ggf.\ \enquote{\textit{Beschriftung im Editor}}). Richtung ist
        \enquote{ein} oder \enquote{aus}, die Seite \enquote{vorn} oder \enquote{hinten} bei
        zweiseitigen Elementen.
  \item[Ports] wie die Licht-Anschlüsse in der Form mit Ports und Verbindungen zu Ports
        zusammengefasst sind: Ein Port bündelt einen Eingang und einen Ausgang derselben Stelle.
  \item[Parameter] Name (darunter ggf.\ die Beschriftung im Dialog), Typ, Standardwert,
        Einheit oder Wertebereich und Bedeutung.
\end{description}

\paragraph{Übersichtsbild.} Der Reiter \gui{Overview} im Dialog \gui{Parameters…} zeichnet
\block{rtt.*} (Raytracer, Abschnitt~\ref{sec:ref-rt-uebersicht}), \block{fg.grating} und
\block{fg.stack} (FourierGrating, Abschnitt~\ref{sec:fg-ref-uebersicht}) sowie \block{optics.*}
und \block{fusion.particle} (Workflow-Plugin) als Zeichnung ihres Teilmodells; alle anderen Blöcke
zeigen dort ihr Schema mit Ports, Anschlüssen und Parametern (Abschnitt~\ref{sec:workflow-uebersicht}).

\paragraph{Darstellungen.} Die Kennbuchstaben der Licht-Anschlüsse bedeuten:
\code{R} Strahlen, \code{H} hybrides Feld, \code{G} Gitterfeld, \code{S} Ebene-Wellen-Spektrum,
\code{T} Jones-Tabelle, \code{P} radiometrisch ohne Phase
(Abschnitt~\ref{sec:workflow-darstellungen}). Gerechnet werden derzeit R und, mit
\code{fieldtracing}, H und S. Welche Blöcke S \emph{nativ} nehmen (Gitter, Stapel, verdichtete
Gruppe, Partikel, Detektoren, \block{fusion.convert}) und wann S lohnt, steht in
Abschnitt~\ref{sec:workflow-darstellungen}; alle übrigen Blöcke mit S am Anschluss bekommen an
ihrem Port ein Feld oder Strahlen.

\paragraph{Datentypen.} \code{result} (ein Ergebnis: Bild, Kurven, Zahlen; der Ausgang
\code{view} eines Anzeigeblocks ist ebenfalls \code{result}), \code{drive} (ein Antriebswert
eines Meta-Blocks), \code{N} (Brechzahlprofil), \code{Q} (Dipolquellen).

\paragraph{Gemeinsame Anschlüsse.} Jeder Block der Kategorien \code{source}, \code{domain},
\code{splitter} und \code{detector} hat zusätzlich den Daten-Eingang \code{pose} (Datentyp
\code{drive}): eine Lageänderung aus sechs Werten (\code{dx}, \code{dy}, \code{dz} in mm,
\code{rx}, \code{ry}, \code{rz} in rad) im lokalen System des Blocks, gedreht um seinen
Bezugspunkt (\param{point} bzw.\ \param{center}). Er wird von einem Meta-Block angetrieben
und ist in den Tabellen nicht wiederholt.

\paragraph{Dynamische Anschlüsse.} Die Raytatouille-Blöcke \block{rtt.section},
\block{rtt.mirror}, \block{rtt.splitter} und \block{rtt.detector} bekommen ihre Anschlüsse und Ports
beim Import aus dem optischen System; die Tabellen zeigen die Grundform. Ihre Daten (Elemente,
Flächen, Gläser, Durchläufe) stehen in internen Parametern und erscheinen im Parameterdialog als
Reiter, die der Raytracer beschreibt.

\paragraph{Typen von Vektoren.} Parameter vom Typ \enquote{Vektor} sind drei Zahlen, im
Parameterdialog durch Kommas getrennt eingegeben (\code{0, 0, 1}); Punkte in mm, Richtungen
ohne Einheit. Parameter vom Typ \enquote{Daten} sind interne Daten, die der
Dialog nicht direkt zeigt.

\section{Übersicht}\label{sec:bloecke-uebersicht}

\begin{referenz}{@{}>{\raggedright\arraybackslash}p{0.3\linewidth}>{\raggedright\arraybackslash}p{0.15\linewidth}>{\raggedright\arraybackslash}p{0.17\linewidth}X@{}}
\toprule
\textbf{Typ} & \textbf{Kategorie} & \textbf{Anbieter} & \textbf{Kurz} \\
\midrule
\endhead
\block{rtt.source} & source & Raytracer & Quelle eines Systems \\
\block{optics.collimated\_source} & source & Workflow & kollimiertes Bündel \\
\block{rtt.section} & domain & Raytracer & Flächenfolge eines Systems \\
\block{rtt.mirror} & domain & Raytracer & Spiegel eines Systems \\
\block{rtt.part.lens} & domain & Raytracer & neue Linse \\
\block{rtt.part.mirror} & domain & Raytracer & neuer Spiegel \\
\block{rtt.part.stop} & domain & Raytracer & neue Blende \\
\block{optics.plane\_mirror} & domain & Workflow & ebener Spiegel \\
\block{optics.interface} & domain & Workflow & ebene Grenzfläche \\
\block{optics.thin\_lens} & domain & Workflow & dünne Linse \\
\block{optics.aplanatic\_lens} & domain & Workflow & ideale aplanatische Linse \\
\block{fg.stack} & domain & FourierGrating & Schichtstapel \\
\block{fusion.particle} & domain & Workflow & Partikel (CliffordBem) \\
\block{rtt.splitter} & splitter & Raytracer & Strahlteiler eines Systems \\
\block{rtt.part.splitter} & splitter & Raytracer & neuer Strahlteiler \\
\block{optics.beam\_splitter} & splitter & Workflow & ebener Strahlteiler \\
\block{optics.splitter} & splitter & Workflow & Teiler mit vier Ports \\
\block{fg.grating} & splitter & FourierGrating & Gitter \\
\block{hpfem.grating} & splitter & hp-FEM & Gitter (hp-FEM) \\
\block{rtt.detector} & detector & Raytracer & Detektor eines Systems \\
\block{rtt.part.detector} & detector & Raytracer & neuer Detektor \\
\block{optics.detector} & detector & Workflow & ebener Detektor \\
\block{fusion.convert} & conversion & Workflow & Umwandlung \\
\block{analysis.camera\_image} & analysis & Workflow & Kamerabild \\
\block{biooptomd.layer} & data & Workflow & BioOptoMD-Daten \\
\block{pathtracer.bsdf\_file} & data & Path Tracer & BSDF aus Datei \\
\block{pathtracer.bsdf\_analytic} & data & Path Tracer & BSDF aus Modell \\
\block{pathtracer.bsdf\_mix} & data & Path Tracer & BSDF-Mischung \\
\block{pathtracer.material} & data & Path Tracer & Material des Path Tracers \\
\block{pathtracer.nk} & data & Path Tracer & Brechzahl \\
\block{pathtracer.topography} & data & Path Tracer & raue Oberfläche \\
\block{pathtracer.rough\_surface} & data & Path Tracer & Streuung rauer Flächen (Näherung) \\
\block{display.image} & analysis & Workflow & Bildanzeige \\
\block{display.plot} & analysis & Workflow & Kurvenanzeige \\
\block{display.beam} & analysis & Workflow & Sonde \\
\block{meta.sweep} & meta & Workflow & Parameterstudie \\
\block{meta.replicates} & meta & Workflow & Wiederholungen \\
\block{meta.optimise} & meta & Workflow & Optimierung \\
\block{group} & domain & Workflow & Gruppe \\
\block{fusion.condensed} & domain & Workflow & verdichtete Gruppe \\
\block{fusion.sum} & splitter & Workflow & Knotenpunkt \\
\bottomrule
\end{referenz}

\section{Quellen}\label{sec:bloecke-quellen}

\subsection{\block{rtt.source}}\label{sec:block-rtt-source}\index{rtt.source@\texttt{rtt.source} (Block)}

\textbf{Kategorie:} \code{source} \quad \textbf{Anbieter:} \plugin{Raytracer} (gerechnet in seinem Prozess; ohne Raytracer rechnet die Workflow-Maschine selbst, mit \code{raytatouille})

\textit{Field points, wavelengths, polarisation of a system (coherence groups)}

Die Lichtquelle eines Raytatouille-Systems: Feldpunkte, Wellenlängen und Polarisation
(Kohärenzgruppen). Sie entsteht mit \gui{Workflow from system} oder \gui{New Raytatouille system}
und hält die Einstellungen des ganzen Systems. Im Parameterdialog zeigt sie die
Systemeinstellungen (Wellenlängen, Felder, Apertur, Objekt) als eigene Felder, die der Raytracer
beschreibt. Die Zahl der Strahlen je Feld und Wellenlänge kommt aus dem Panel des
Workflow-Plugins (\gui{Rays per field and wavelength}).

\textbf{Licht:} \code{out} (aus, RHGSP, \enquote{light}).

\textbf{Ports:} \code{out} (aus \code{out}).

\begin{referenz}{@{}>{\raggedright\arraybackslash}p{0.2\linewidth}>{\raggedright\arraybackslash}p{0.11\linewidth}>{\raggedright\arraybackslash}p{0.13\linewidth}>{\raggedright\arraybackslash}p{0.15\linewidth}>{\raggedright\arraybackslash}X@{}}
\toprule
\textbf{Parameter} & \textbf{Typ} & \textbf{Standard} & \textbf{Einheit / Bereich} & \textbf{Bedeutung} \\
\midrule
\endhead
\param{system} \newline\small\gui{Raytatouille system} & Datei & leer & -- & Systemdatei (\datei{.rtt.json}) eines älteren Workflows; neue Workflows halten das System in den Daten der Blöcke. \\
\param{fields} & Text & leer & -- & Indizes der Feldpunkte, die gerechnet werden; leer: alle. \\
\param{wavelengths} & Text & leer & -- & Indizes der Wellenlängen; leer: alle. \\
\param{rtt\_\allowbreak{}system} & Daten & -- & -- & Einstellungen des Systems (Wellenlängen, Felder, Apertur, Objekt); im Dialog als eigene Felder. \\
\param{rtt\_\allowbreak{}tree} & Daten & -- & -- & Knoten des Raytatouille-Modellbaums, die kein Block besitzt. \\
\param{rtt\_\allowbreak{}paths} & Daten & -- & -- & Reihenfolge und Schlüssel der Pfade des Systems. \\
\param{rtt\_\allowbreak{}catalogs} & Daten & -- & -- & Glaskatalog-Dateien (vom Raytracer-Panel übernommen). \\
\param{rtt\_\allowbreak{}catalog\_\allowbreak{}data} & Daten & -- & -- & Inhalt der Glaskataloge, im Dokument gespeichert. \\
\bottomrule
\end{referenz}

\subsection{\block{optics.collimated\_source}}\label{sec:block-optics-collimated-source}\index{optics.collimatedsource@\texttt{optics.collimated\_source} (Block)}

\textbf{Kategorie:} \code{source} \quad \textbf{Anbieter:} \plugin{Workflow} (RTT-Fusion, eingebaut)

\textit{Collimated beam (plane wave per field angle and wavelength)}

Ein kollimiertes Bündel ohne Raytatouille: je Feldwinkel und Wellenlänge eine ebene Welle über
einer Kreisscheibe, auf einem hexapolaren Raster abgetastet. Für einfache Interferometer
(Michelson, Linnik ohne Objektive) und Tests.

Die Polarisation ist linear längs \param{polarisation} oder – mit \param{ellipticity\_deg} – elliptisch
bzw.\ zirkular; die Strahlen tragen sie als komplexen Feldvektor durch alle Blöcke und
Umwandlungen. Beispiel: \param{direction} \code{0, 0, 1}, \param{polarisation} \code{1, 0, 0},
\param{ellipticity\_deg} \code{45} ergibt das Feld $(1, \mathrm{i}, 0)/\sqrt 2$, das von $x$ nach
$y$ dreht (Zeitabhängigkeit $e^{-\mathrm{i}\omega t}$); \code{-45} dreht andersherum.

\textbf{Licht:} \code{out} (aus, RHS).

\textbf{Ports:} \code{out} (aus \code{out}).

\begin{referenz}{@{}>{\raggedright\arraybackslash}p{0.2\linewidth}>{\raggedright\arraybackslash}p{0.11\linewidth}>{\raggedright\arraybackslash}p{0.13\linewidth}>{\raggedright\arraybackslash}p{0.15\linewidth}>{\raggedright\arraybackslash}X@{}}
\toprule
\textbf{Parameter} & \textbf{Typ} & \textbf{Standard} & \textbf{Einheit / Bereich} & \textbf{Bedeutung} \\
\midrule
\endhead
\param{center} & Vektor & \code{0.0, 0.0, 0.0} & mm & Mittelpunkt des Bündels. \\
\param{direction} & Vektor & \code{0.0, 0.0, 1.0} & -- & Ausbreitungsrichtung (Achse) des Bündels. \\
\param{radius} & Zahl & \code{1.0} & mm; $\geq$ 0.0 & Radius des Bündels. \\
\param{rings} & Ganzzahl & \code{20} & $\geq$ 1 & Ringe des hexapolaren Strahlrasters ($1 + 3r(r+1)$ Strahlen). \\
\param{wavelengths\_\allowbreak{}um} & Zahlenliste & \code{0.6328} & \si{\micro\metre} & Wellenlängen. \\
\param{field\_\allowbreak{}angles\_\allowbreak{}deg} & Zahlenliste & \code{(0.0, 0.0)} & \si{\degree} & Feldwinkel als Paare, z.\,B.\ \code{[0, 0], [1, 0]}. \\
\param{irradiance} & Zahl & \code{1.0} & W/mm$^2$ & Bestrahlungsstärke (bei Gauß-Profil: auf der Achse). \\
\param{polarisation} & Vektor & \code{1.0, 0.0, 0.0} & -- & Richtung des elektrischen Felds; bei elliptischem Licht die große Halbachse $\hat p$. \\
\param{ellipticity\_\allowbreak{}deg} & Zahl & \code{0.0} & \si{\degree}; $-45$ \dots\ 45 & Elliptizität $\chi$: 0 linear, $\pm 45$ zirkular, $\tan\chi$ = kleine/große Halbachse. Das Feld ist $\cos\chi\,\hat p + \mathrm{i}\sin\chi\,(\hat d \times \hat p)$; bei $\chi > 0$ dreht es von $\hat p$ nach $\hat d \times \hat p$ ($S_3 > 0$). \\
\param{up} & Vektor & \code{0.0, 1.0, 0.0} & -- & Richtung der lokalen v-Achse. \\
\param{gaussian\_\allowbreak{}waist} & Zahl & \code{0.0} & mm; $\geq$ 0.0 & 0: gleichmäßig; $> 0$: Gauß-Amplitude $\exp(-\rho^2/w^2)$ mit diesem $w$. \\
\bottomrule
\end{referenz}

\section{Optische Elemente (Domänen)}\label{sec:bloecke-elemente}

\subsection{\block{rtt.section}}\label{sec:block-rtt-section}\index{rtt.section@\texttt{rtt.section} (Block)}

\textbf{Kategorie:} \code{domain} \quad \textbf{Anbieter:} \plugin{Raytracer} (gerechnet in seinem Prozess; ohne Raytracer rechnet die Workflow-Maschine selbst, mit \code{raytatouille})

\textit{Sequential surfaces of a Raytatouille system (e.g. an objective)}

Sequentielle Flächen eines Raytatouille-Systems, z.\,B.\ ein Objektiv. Licht durchläuft den
Block in beiden Richtungen (vorn hinein und hinten heraus oder umgekehrt). Die Anschlüsse
werden beim Import aus dem System erzeugt; der Dialog zeigt einen Reiter je Element mit Radien,
Kegelschnittkonstanten, Aperturen, Flächenlagen und Gläsern. Mit \gui{Section $\to$ group} wird
daraus eine Gruppe mit einem Block je Element. Der Daten-Eingang \code{drive} verschiebt den
Block längs der Flächennormalen (Phasenschieber).

\textbf{Licht:} \code{front\_in} (ein, RHGSTP, vorn, \enquote{front in}); \code{back\_out} (aus, RHGSTP, hinten, \enquote{back out}); \code{back\_in} (ein, RHGSTP, hinten, \enquote{back in}); \code{front\_out} (aus, RHGSTP, vorn, \enquote{front out}).\\
\textbf{Daten:} \code{drive} (ein, \code{drive}).

\textbf{Ports:} \code{front} (ein \code{front\_in}, aus \code{front\_out}); \code{back} (ein \code{back\_in}, aus \code{back\_out}).

\begin{referenz}{@{}>{\raggedright\arraybackslash}p{0.2\linewidth}>{\raggedright\arraybackslash}p{0.11\linewidth}>{\raggedright\arraybackslash}p{0.13\linewidth}>{\raggedright\arraybackslash}p{0.15\linewidth}>{\raggedright\arraybackslash}X@{}}
\toprule
\textbf{Parameter} & \textbf{Typ} & \textbf{Standard} & \textbf{Einheit / Bereich} & \textbf{Bedeutung} \\
\midrule
\endhead
\param{location} & Text & leer & -- & Knoten des Raytatouille-Modellbaums, den der Block darstellt. \\
\param{offset} & Vektor & \code{0.0, 0.0, 0.0} & mm & Verschiebung der Elemente des Blocks. \\
\param{drive} & Zahl & \code{0.0} & mm & Verschiebung längs der Flächennormalen (vom \code{drive} eines Phasenschiebers gesetzt). \\
\bottomrule
\end{referenz}

\subsection{\block{rtt.mirror}}\label{sec:block-rtt-mirror}\index{rtt.mirror@\texttt{rtt.mirror} (Block)}

\textbf{Kategorie:} \code{domain} \quad \textbf{Anbieter:} \plugin{Raytracer} (gerechnet in seinem Prozess; ohne Raytracer rechnet die Workflow-Maschine selbst, mit \code{raytatouille})

\textit{Mirror (one surface, reflecting)}

Ein Spiegel eines Raytatouille-Systems (eine reflektierende Fläche). Licht kommt durch den Port
\code{front} herein und verlässt ihn durch denselben Port.

\textbf{Licht:} \code{in} (ein, RHGS); \code{out} (aus, RHGS).\\
\textbf{Daten:} \code{drive} (ein, \code{drive}).

\textbf{Ports:} \code{front} (ein \code{in}, aus \code{out}).

\begin{referenz}{@{}>{\raggedright\arraybackslash}p{0.2\linewidth}>{\raggedright\arraybackslash}p{0.11\linewidth}>{\raggedright\arraybackslash}p{0.13\linewidth}>{\raggedright\arraybackslash}p{0.15\linewidth}>{\raggedright\arraybackslash}X@{}}
\toprule
\textbf{Parameter} & \textbf{Typ} & \textbf{Standard} & \textbf{Einheit / Bereich} & \textbf{Bedeutung} \\
\midrule
\endhead
\param{surface} & Text & leer & -- & Kennung der Fläche im Raytatouille-System. \\
\param{offset} & Vektor & \code{0.0, 0.0, 0.0} & mm & Verschiebung der Elemente des Blocks. \\
\param{drive} & Zahl & \code{0.0} & mm & Verschiebung längs der Flächennormalen (vom \code{drive} eines Phasenschiebers gesetzt). \\
\bottomrule
\end{referenz}

\subsection{\block{rtt.part.lens}}\label{sec:block-rtt-part-lens}\index{rtt.part.lens@\texttt{rtt.part.lens} (Block)}

\textbf{Kategorie:} \code{domain} \quad \textbf{Anbieter:} \plugin{Raytracer} (gerechnet in seinem Prozess; ohne Raytracer rechnet die Workflow-Maschine selbst, mit \code{raytatouille})

\textit{Raytatouille lens (a new element)}

Eine neue Raytatouille-Linse aus der Bibliothek (nicht importiert). Ihr Element – Lage, Flächen
mit Form und Apertur, Glas – steht im Parameter \param{rtt\_new} und wird im Parameterdialog wie
die Daten eines importierten Blocks bearbeitet. Vorlage: Linse aus \code{CONST:1.5168} mit den
Radien 51{,}68~mm und $-51{,}68$~mm, 4~mm dick, kreisförmige Apertur mit 12{,}5~mm Radius.
Verlässt das Licht den Block durch die Seite, durch die es kam, ist das eine Reflexion, sonst
eine Transmission.

\textbf{Licht:} \code{front\_in} (ein, RHS, vorn); \code{front\_out} (aus, RHS, vorn); \code{back\_in} (ein, RHS, hinten); \code{back\_out} (aus, RHS, hinten).

\textbf{Ports:} \code{front} (ein \code{front\_in}, aus \code{front\_out}); \code{back} (ein \code{back\_in}, aus \code{back\_out}).

\begin{referenz}{@{}>{\raggedright\arraybackslash}p{0.2\linewidth}>{\raggedright\arraybackslash}p{0.11\linewidth}>{\raggedright\arraybackslash}p{0.13\linewidth}>{\raggedright\arraybackslash}p{0.15\linewidth}>{\raggedright\arraybackslash}X@{}}
\toprule
\textbf{Parameter} & \textbf{Typ} & \textbf{Standard} & \textbf{Einheit / Bereich} & \textbf{Bedeutung} \\
\midrule
\endhead
\param{rtt\_\allowbreak{}new} & Daten & -- & -- & Das Raytatouille-Element (Lage, Flächen, Glas); im Dialog als eigene Felder. \\
\bottomrule
\end{referenz}

\subsection{\block{rtt.part.mirror}}\label{sec:block-rtt-part-mirror}\index{rtt.part.mirror@\texttt{rtt.part.mirror} (Block)}

\textbf{Kategorie:} \code{domain} \quad \textbf{Anbieter:} \plugin{Raytracer} (gerechnet in seinem Prozess; ohne Raytracer rechnet die Workflow-Maschine selbst, mit \code{raytatouille})

\textit{Raytatouille mirror (a new element)}

Ein neuer Raytatouille-Spiegel aus der Bibliothek. Vorlage: idealer Spiegel, kreisförmige
Apertur mit 12{,}5~mm Radius.

\textbf{Licht:} \code{front\_in} (ein, RHS, vorn); \code{front\_out} (aus, RHS, vorn).

\textbf{Ports:} \code{front} (ein \code{front\_in}, aus \code{front\_out}).

\begin{referenz}{@{}>{\raggedright\arraybackslash}p{0.2\linewidth}>{\raggedright\arraybackslash}p{0.11\linewidth}>{\raggedright\arraybackslash}p{0.13\linewidth}>{\raggedright\arraybackslash}p{0.15\linewidth}>{\raggedright\arraybackslash}X@{}}
\toprule
\textbf{Parameter} & \textbf{Typ} & \textbf{Standard} & \textbf{Einheit / Bereich} & \textbf{Bedeutung} \\
\midrule
\endhead
\param{rtt\_\allowbreak{}new} & Daten & -- & -- & Das Raytatouille-Element (Lage, Flächen, Glas); im Dialog als eigene Felder. \\
\bottomrule
\end{referenz}

\subsection{\block{rtt.part.stop}}\label{sec:block-rtt-part-stop}\index{rtt.part.stop@\texttt{rtt.part.stop} (Block)}

\textbf{Kategorie:} \code{domain} \quad \textbf{Anbieter:} \plugin{Raytracer} (gerechnet in seinem Prozess; ohne Raytracer rechnet die Workflow-Maschine selbst, mit \code{raytatouille})

\textit{Raytatouille stop (a new element)}

Eine neue Raytatouille-Blende aus der Bibliothek. Vorlage: kreisförmige Öffnung mit 5~mm Radius.
\gui{New Raytatouille system} legt eine solche Blende bei $z = 0$ an.

\textbf{Licht:} \code{front\_in} (ein, RHS, vorn); \code{front\_out} (aus, RHS, vorn); \code{back\_in} (ein, RHS, hinten); \code{back\_out} (aus, RHS, hinten).

\textbf{Ports:} \code{front} (ein \code{front\_in}, aus \code{front\_out}); \code{back} (ein \code{back\_in}, aus \code{back\_out}).

\begin{referenz}{@{}>{\raggedright\arraybackslash}p{0.2\linewidth}>{\raggedright\arraybackslash}p{0.11\linewidth}>{\raggedright\arraybackslash}p{0.13\linewidth}>{\raggedright\arraybackslash}p{0.15\linewidth}>{\raggedright\arraybackslash}X@{}}
\toprule
\textbf{Parameter} & \textbf{Typ} & \textbf{Standard} & \textbf{Einheit / Bereich} & \textbf{Bedeutung} \\
\midrule
\endhead
\param{rtt\_\allowbreak{}new} & Daten & -- & -- & Das Raytatouille-Element (Lage, Flächen, Glas); im Dialog als eigene Felder. \\
\bottomrule
\end{referenz}

\subsection{\block{optics.plane\_mirror}}\label{sec:block-optics-plane-mirror}\index{optics.planemirror@\texttt{optics.plane\_mirror} (Block)}

\textbf{Kategorie:} \code{domain} \quad \textbf{Anbieter:} \plugin{Workflow} (RTT-Fusion, eingebaut)

\textit{Plane mirror}

Ein ebener Spiegel (analytisch, ohne Raytatouille). Die Ebene geht durch \param{point} mit der
Normalen \param{normal}; zwischen den Elementen breitet sich das Licht frei in Luft ($n = 1$) aus.

\textbf{Licht:} \code{in} (ein, RHS); \code{out} (aus, RHS).\\
\textbf{Daten:} \code{drive} (ein, \code{drive}).

\textbf{Ports:} \code{front} (ein \code{in}, aus \code{out}).

\begin{referenz}{@{}>{\raggedright\arraybackslash}p{0.2\linewidth}>{\raggedright\arraybackslash}p{0.11\linewidth}>{\raggedright\arraybackslash}p{0.13\linewidth}>{\raggedright\arraybackslash}p{0.15\linewidth}>{\raggedright\arraybackslash}X@{}}
\toprule
\textbf{Parameter} & \textbf{Typ} & \textbf{Standard} & \textbf{Einheit / Bereich} & \textbf{Bedeutung} \\
\midrule
\endhead
\param{point} & Vektor & \code{0.0, 0.0, 0.0} & mm & Bezugspunkt der Ebene bzw.\ des Elements. \\
\param{normal} & Vektor & \code{0.0, 0.0, -1.0} & -- & Flächennormale. \\
\param{up} & Vektor & \code{0.0, 1.0, 0.0} & -- & Richtung der lokalen v-Achse. \\
\param{drive} & Zahl & \code{0.0} & mm & Verschiebung längs der Normalen. \\
\param{reflectance\_\allowbreak{}amplitude} & Zahl & \code{-1.0} & -- & Amplitudenreflexionsgrad (Standard $-1$: idealer Spiegel mit Phasensprung). \\
\bottomrule
\end{referenz}

\subsection{\block{optics.interface}}\label{sec:block-optics-interface}\index{optics.interface@\texttt{optics.interface} (Block)}

\textbf{Kategorie:} \code{domain} \quad \textbf{Anbieter:} \plugin{Workflow} (RTT-Fusion, eingebaut)

\textit{Plane interface of two media (Fresnel, s and p)}

Eine ebene Grenzfläche zweier Medien mit Fresnel-Koeffizienten für s- und p-Polarisation,
in beiden Richtungen durchlaufbar.

\textbf{Licht:} \code{front\_in} (ein, RHS, vorn); \code{front\_out} (aus, RHS, vorn); \code{back\_in} (ein, RHS, hinten); \code{back\_out} (aus, RHS, hinten).

\textbf{Ports:} \code{front} (ein \code{front\_in}, aus \code{front\_out}); \code{back} (ein \code{back\_in}, aus \code{back\_out}).

\begin{referenz}{@{}>{\raggedright\arraybackslash}p{0.2\linewidth}>{\raggedright\arraybackslash}p{0.11\linewidth}>{\raggedright\arraybackslash}p{0.13\linewidth}>{\raggedright\arraybackslash}p{0.15\linewidth}>{\raggedright\arraybackslash}X@{}}
\toprule
\textbf{Parameter} & \textbf{Typ} & \textbf{Standard} & \textbf{Einheit / Bereich} & \textbf{Bedeutung} \\
\midrule
\endhead
\param{point} & Vektor & \code{0.0, 0.0, 0.0} & mm & Bezugspunkt der Ebene bzw.\ des Elements. \\
\param{normal} & Vektor & \code{0.0, 0.0, -1.0} & -- & Flächennormale. \\
\param{up} & Vektor & \code{0.0, 1.0, 0.0} & -- & Richtung der lokalen v-Achse. \\
\param{drive} & Zahl & \code{0.0} & mm & Verschiebung längs der Normalen. \\
\param{n\_\allowbreak{}front} & Zahl & \code{1.0} & -- & Brechzahl auf der Seite, in die die Normale zeigt. \\
\param{n\_\allowbreak{}back} & Zahl & \code{1.5} & -- & Brechzahl dahinter. \\
\bottomrule
\end{referenz}

\subsection{\block{optics.thin\_lens}}\label{sec:block-optics-thin-lens}\index{optics.thinlens@\texttt{optics.thin\_lens} (Block)}

\textbf{Kategorie:} \code{domain} \quad \textbf{Anbieter:} \plugin{Workflow} (RTT-Fusion, eingebaut)

\textit{Paraxial thin lens (matrix optics, any incidence)}

Eine paraxiale dünne Linse (Matrixoptik, beliebiger Einfall).

\textbf{Licht:} \code{in} (ein, RHS); \code{out} (aus, RHS).

\textbf{Ports:} \code{in} (ein \code{in}); \code{out} (aus \code{out}).

\begin{referenz}{@{}>{\raggedright\arraybackslash}p{0.2\linewidth}>{\raggedright\arraybackslash}p{0.11\linewidth}>{\raggedright\arraybackslash}p{0.13\linewidth}>{\raggedright\arraybackslash}p{0.15\linewidth}>{\raggedright\arraybackslash}X@{}}
\toprule
\textbf{Parameter} & \textbf{Typ} & \textbf{Standard} & \textbf{Einheit / Bereich} & \textbf{Bedeutung} \\
\midrule
\endhead
\param{point} & Vektor & \code{0.0, 0.0, 0.0} & mm & Bezugspunkt der Ebene bzw.\ des Elements. \\
\param{axis} & Vektor & \code{0.0, 0.0, 1.0} & -- & Optische Achse. \\
\param{focal\_\allowbreak{}length} & Zahl & \code{10.0} & mm & Brennweite. \\
\param{aperture\_\allowbreak{}radius} & Zahl & $10^{9}$ & mm & Radius der Apertur. \\
\bottomrule
\end{referenz}

\subsection{\block{optics.aplanatic\_lens}}\label{sec:block-optics-aplanatic-lens}\index{optics.aplanaticlens@\texttt{optics.aplanatic\_lens} (Block)}

\textbf{Kategorie:} \code{domain} \quad \textbf{Anbieter:} \plugin{Workflow} (RTT-Fusion, eingebaut)

\textit{Ideal aplanatic lens (sine condition, Richards-Wolf illumination)}

Eine ideale aplanatische Linse (Sinusbedingung, Beleuchtung nach Richards-Wolf), z.\,B.\ als
Ersatz eines realen Objektivs in einer Variante einer Gruppe.

\textbf{Licht:} \code{in} (ein, RHS); \code{out} (aus, RHS).

\textbf{Ports:} \code{in} (ein \code{in}); \code{out} (aus \code{out}).

\begin{referenz}{@{}>{\raggedright\arraybackslash}p{0.2\linewidth}>{\raggedright\arraybackslash}p{0.11\linewidth}>{\raggedright\arraybackslash}p{0.13\linewidth}>{\raggedright\arraybackslash}p{0.15\linewidth}>{\raggedright\arraybackslash}X@{}}
\toprule
\textbf{Parameter} & \textbf{Typ} & \textbf{Standard} & \textbf{Einheit / Bereich} & \textbf{Bedeutung} \\
\midrule
\endhead
\param{point} & Vektor & \code{0.0, 0.0, 0.0} & mm & Bezugspunkt der Ebene bzw.\ des Elements. \\
\param{axis} & Vektor & \code{0.0, 0.0, 1.0} & -- & Optische Achse. \\
\param{focal\_\allowbreak{}length} & Zahl & \code{10.0} & mm & Brennweite. \\
\param{aperture\_\allowbreak{}radius} & Zahl & \code{5.0} & mm & Radius der Apertur. \\
\bottomrule
\end{referenz}

\subsection{\block{fg.stack}}\label{sec:block-fg-stack}\index{fg.stack@\texttt{fg.stack} (Block)}

\textbf{Kategorie:} \code{domain} \quad \textbf{Anbieter:} \plugin{FourierGrating} (gerechnet in seinem Prozess)

\textit{Layer stack (FourierGrating, exact per ray): coating, filter, mirror layers}

Ein Schichtstapel (Beschichtung, Filter, Spiegelschichten) zwischen zwei Medien, exakt je Strahl
mit FourierGrating gerechnet. Die Schichten stehen in einer FourierGrating-Konfiguration
(homogene Schichten vom Substrat hinten bis zum Superstrat vorn). \gui{Check workflow} speichert
die Konfiguration im Block. Mit \gui{Store panel configuration in block} übernehmen Sie die
Konfiguration aus dem FourierGrating-Panel (Kapitel~\ref{kap:fouriergrating}).

\textbf{Licht:} \code{front\_in} (ein, RHS, vorn); \code{front\_out} (aus, RHS, vorn); \code{back\_in} (ein, RHS, hinten); \code{back\_out} (aus, RHS, hinten).

\textbf{Ports:} \code{front} (ein \code{front\_in}, aus \code{front\_out}); \code{back} (ein \code{back\_in}, aus \code{back\_out}).

\begin{referenz}{@{}>{\raggedright\arraybackslash}p{0.2\linewidth}>{\raggedright\arraybackslash}p{0.11\linewidth}>{\raggedright\arraybackslash}p{0.13\linewidth}>{\raggedright\arraybackslash}p{0.15\linewidth}>{\raggedright\arraybackslash}X@{}}
\toprule
\textbf{Parameter} & \textbf{Typ} & \textbf{Standard} & \textbf{Einheit / Bereich} & \textbf{Bedeutung} \\
\midrule
\endhead
\param{fouriergrating} & Datei & leer & -- & FourierGrating-Konfiguration des Stapels; \gui{Check workflow} speichert sie im Block (\param{config}) und leert den Pfad. \\
\param{config} & Daten & -- & -- & Die im Block gehaltene Konfiguration (hat Vorrang vor der Datei). \\
\param{point} & Vektor & \code{0.0, 0.0, 0.0} & mm & Bezugspunkt des Stapels. \\
\param{normal} & Vektor & \code{0.0, 0.0, 1.0} & -- & Normale, zeigt nach vorn (zum Superstrat). \\
\param{field\_\allowbreak{}threshold} & Zahl & $10^{-10}$ & -- & Ebene Wellen unter diesem Anteil der größten Amplitude werden weggelassen (H). \\
\param{field\_\allowbreak{}transfer} & Auswahl & \code{exact} & \code{exact}, \code{interpolated} & \code{exact}: FourierGrating je ebener Welle; \code{interpolated}: aus einem geprüften Teilraster des Spektrums ($10^{-8}$), viel schneller für breite Bündel. \\
\param{field\_\allowbreak{}grid\_\allowbreak{}max} & Ganzzahl & \code{1024} & -- & Höchstens so viele Gitterpunkte je Achse für ein Feld (H). \\
\bottomrule
\end{referenz}

\subsection{\block{fusion.particle}}\label{sec:block-fusion-particle}\index{fusion.particle@\texttt{fusion.particle} (Block)}

\textbf{Kategorie:} \code{domain} \quad \textbf{Anbieter:} \plugin{Workflow} (RTT-Fusion, eingebaut); braucht das Paket \code{cliffordbem}

\textit{Particle (CliffordBem BEM): backscattered light of a focused beam}

Ein kugelförmiges Partikel als Probe, gerechnet mit dem BEM-Löser CliffordBem: Die einfallenden
Strahlen (z.\,B.\ der Fokus eines Objektivs) werden in ein Winkelspektrum umgerechnet, am
Partikel gestreut und als rückgestreute Strahlen zurückgeschickt. Gleiche Anschlüsse wie ein
Spiegel, sodass er einen Probenspiegel mit \gui{Replace with} ersetzen kann.

\textbf{Licht:} \code{in} (ein, RHS); \code{out} (aus, RHS, \enquote{backscattered}).

\textbf{Ports:} \code{front} (ein \code{in}, aus \code{out}).

\begin{referenz}{@{}>{\raggedright\arraybackslash}p{0.2\linewidth}>{\raggedright\arraybackslash}p{0.11\linewidth}>{\raggedright\arraybackslash}p{0.13\linewidth}>{\raggedright\arraybackslash}p{0.15\linewidth}>{\raggedright\arraybackslash}X@{}}
\toprule
\textbf{Parameter} & \textbf{Typ} & \textbf{Standard} & \textbf{Einheit / Bereich} & \textbf{Bedeutung} \\
\midrule
\endhead
\param{point} & Vektor & \code{0.0, 0.0, 0.0} & mm & Mittelpunkt des Partikels. \\
\param{radius\_\allowbreak{}um} & Zahl & \code{0.25} & \si{\micro\metre}; $\geq$ 0.0 & Radius der Kugel. \\
\param{index} & Zahl & \code{1.5} & -- & Brechzahl des Partikels. \\
\param{n\_\allowbreak{}medium} & Zahl & \code{1.0} & -- & Brechzahl des umgebenden Mediums. \\
\param{refinement} & Ganzzahl & \code{4} & 1 … 7 & Verfeinerungsstufe des Ikosaeder-Netzes. \\
\param{n\_\allowbreak{}theta} & Ganzzahl & \code{40} & $\geq$ 4 & Stützstellen der Fernfeld-Quadratur in $\theta$. \\
\param{n\_\allowbreak{}phi} & Ganzzahl & \code{80} & $\geq$ 8 & Stützstellen der Fernfeld-Quadratur in $\varphi$. \\
\bottomrule
\end{referenz}

\section{Strahlteiler und Gitter}\label{sec:bloecke-teiler}

\subsection{\block{rtt.splitter}}\label{sec:block-rtt-splitter}\index{rtt.splitter@\texttt{rtt.splitter} (Block)}

\textbf{Kategorie:} \code{splitter} \quad \textbf{Anbieter:} \plugin{Raytracer} (gerechnet in seinem Prozess; ohne Raytracer rechnet die Workflow-Maschine selbst, mit \code{raytatouille})

\textit{Beam splitter: splits into channels, combines coherently per group}

Ein Strahlteiler eines Raytatouille-Systems. Licht aus \code{in} wird reflektiert (\code{r}) und
durchgelassen (\code{t}); zurückkehrendes Licht auf \code{r\_in} und \code{t\_in} wird kohärent je
Kohärenzgruppe zu \code{merged} vereinigt. Der Daten-Eingang \code{drive} verschiebt den Teiler
längs seiner Normalen.

\textbf{Licht:} \code{in} (ein, RHGS); \code{r} (aus, RHGS, \enquote{reflected}); \code{t} (aus, RHGS, \enquote{transmitted}); \code{r\_in} (ein, RHGS); \code{t\_in} (ein, RHGS); \code{merged} (aus, RHGS, \enquote{combined}).\\
\textbf{Daten:} \code{drive} (ein, \code{drive}).

\textbf{Ports:} \code{in} (ein \code{in}); \code{r} (ein \code{r\_in}, aus \code{r}); \code{t} (ein \code{t\_in}, aus \code{t}); \code{merged} (aus \code{merged}).

\begin{referenz}{@{}>{\raggedright\arraybackslash}p{0.2\linewidth}>{\raggedright\arraybackslash}p{0.11\linewidth}>{\raggedright\arraybackslash}p{0.13\linewidth}>{\raggedright\arraybackslash}p{0.15\linewidth}>{\raggedright\arraybackslash}X@{}}
\toprule
\textbf{Parameter} & \textbf{Typ} & \textbf{Standard} & \textbf{Einheit / Bereich} & \textbf{Bedeutung} \\
\midrule
\endhead
\param{surface} & Text & leer & -- & Kennung der Fläche im Raytatouille-System. \\
\param{offset} & Vektor & \code{0.0, 0.0, 0.0} & mm & Verschiebung der Elemente des Blocks. \\
\param{drive} & Zahl & \code{0.0} & mm & Verschiebung längs der Flächennormalen (vom \code{drive} eines Phasenschiebers gesetzt). \\
\bottomrule
\end{referenz}

\subsection{\block{rtt.part.splitter}}\label{sec:block-rtt-part-splitter}\index{rtt.part.splitter@\texttt{rtt.part.splitter} (Block)}

\textbf{Kategorie:} \code{splitter} \quad \textbf{Anbieter:} \plugin{Raytracer} (gerechnet in seinem Prozess; ohne Raytracer rechnet die Workflow-Maschine selbst, mit \code{raytatouille})

\textit{Raytatouille beam splitter (a new element)}

Ein neuer Raytatouille-Strahlteiler aus der Bibliothek mit vier Ports. Vorlage: idealer
Strahlteiler, um 45\si{\degree} um x gedreht, Apertur 12{,}5~mm. Verlässt das Licht den Teiler
durch die Seite, durch die es kam, ist das eine Reflexion.

\textbf{Licht:} \code{front\_a\_in} (ein, RHS, vorn); \code{front\_a\_out} (aus, RHS, vorn); \code{front\_b\_in} (ein, RHS, vorn); \code{front\_b\_out} (aus, RHS, vorn); \code{back\_a\_in} (ein, RHS, hinten); \code{back\_a\_out} (aus, RHS, hinten); \code{back\_b\_in} (ein, RHS, hinten); \code{back\_b\_out} (aus, RHS, hinten).

\textbf{Ports:} \code{front\_a} (ein \code{front\_a\_in}, aus \code{front\_a\_out}); \code{front\_b} (ein \code{front\_b\_in}, aus \code{front\_b\_out}); \code{back\_a} (ein \code{back\_a\_in}, aus \code{back\_a\_out}); \code{back\_b} (ein \code{back\_b\_in}, aus \code{back\_b\_out}).

\begin{referenz}{@{}>{\raggedright\arraybackslash}p{0.2\linewidth}>{\raggedright\arraybackslash}p{0.11\linewidth}>{\raggedright\arraybackslash}p{0.13\linewidth}>{\raggedright\arraybackslash}p{0.15\linewidth}>{\raggedright\arraybackslash}X@{}}
\toprule
\textbf{Parameter} & \textbf{Typ} & \textbf{Standard} & \textbf{Einheit / Bereich} & \textbf{Bedeutung} \\
\midrule
\endhead
\param{rtt\_\allowbreak{}new} & Daten & -- & -- & Das Raytatouille-Element (Lage, Flächen, Glas); im Dialog als eigene Felder. \\
\bottomrule
\end{referenz}

\subsection{\block{optics.beam\_splitter}}\label{sec:block-optics-beam-splitter}\index{optics.beamsplitter@\texttt{optics.beam\_splitter} (Block)}

\textbf{Kategorie:} \code{splitter} \quad \textbf{Anbieter:} \plugin{Workflow} (RTT-Fusion, eingebaut)

\textit{Plane beam splitter (Michelson / Linnik arrangement)}

Ein ebener Strahlteiler in Michelson-/Linnik-Anordnung (analytisch): Licht aus \code{in} geht
reflektiert nach \code{r} und durchgelassen nach \code{t}; was auf dem durchgelassenen Arm
zurückkommt (\code{t\_in}), wird nach \code{merged} reflektiert, was auf dem reflektierten Arm
zurückkommt (\code{r\_in}), nach \code{merged} durchgelassen.

\textbf{Licht:} \code{in} (ein, RHS); \code{r} (aus, RHS); \code{t} (aus, RHS); \code{r\_in} (ein, RHS); \code{t\_in} (ein, RHS); \code{merged} (aus, RHS).

\textbf{Ports:} \code{in} (ein \code{in}); \code{r} (ein \code{r\_in}, aus \code{r}); \code{t} (ein \code{t\_in}, aus \code{t}); \code{merged} (aus \code{merged}).

\begin{referenz}{@{}>{\raggedright\arraybackslash}p{0.2\linewidth}>{\raggedright\arraybackslash}p{0.11\linewidth}>{\raggedright\arraybackslash}p{0.13\linewidth}>{\raggedright\arraybackslash}p{0.15\linewidth}>{\raggedright\arraybackslash}X@{}}
\toprule
\textbf{Parameter} & \textbf{Typ} & \textbf{Standard} & \textbf{Einheit / Bereich} & \textbf{Bedeutung} \\
\midrule
\endhead
\param{point} & Vektor & \code{0.0, 0.0, 0.0} & mm & Bezugspunkt der Ebene bzw.\ des Elements. \\
\param{normal} & Vektor & \code{0.0, 0.0, -1.0} & -- & Flächennormale. \\
\param{up} & Vektor & \code{0.0, 1.0, 0.0} & -- & Richtung der lokalen v-Achse. \\
\param{drive} & Zahl & \code{0.0} & mm & Verschiebung längs der Normalen. \\
\param{r} & Zahl & $1/\sqrt{2}$ & -- & Amplitudenreflexionsgrad. \\
\param{t} & Zahl & $1/\sqrt{2}$ & -- & Amplitudentransmissionsgrad. \\
\bottomrule
\end{referenz}

\subsection{\block{optics.splitter}}\label{sec:block-optics-splitter}\index{optics.splitter@\texttt{optics.splitter} (Block)}

\textbf{Kategorie:} \code{splitter} \quad \textbf{Anbieter:} \plugin{Workflow} (RTT-Fusion, eingebaut)

\textit{Lossless splitter with four ports (Stokes relations)}

Ein verlustfreier Teiler mit vier Ports nach den Stokes-Relationen: Licht in einen Port verlässt
den Block durch den anderen Port derselben Seite (reflektiert, z.\,B.\ \code{front\_a} $\to$
\code{front\_b}) und durch den Port mit demselben Buchstaben auf der anderen Seite (durchgelassen,
z.\,B.\ \code{front\_a} $\to$ \code{back\_a}).

\textbf{Licht:} \code{front\_a\_in} (ein, RHS, vorn); \code{front\_a\_out} (aus, RHS, vorn); \code{front\_b\_in} (ein, RHS, vorn); \code{front\_b\_out} (aus, RHS, vorn); \code{back\_a\_in} (ein, RHS, hinten); \code{back\_a\_out} (aus, RHS, hinten); \code{back\_b\_in} (ein, RHS, hinten); \code{back\_b\_out} (aus, RHS, hinten).

\textbf{Ports:} \code{front\_a} (ein \code{front\_a\_in}, aus \code{front\_a\_out}); \code{front\_b} (ein \code{front\_b\_in}, aus \code{front\_b\_out}); \code{back\_a} (ein \code{back\_a\_in}, aus \code{back\_a\_out}); \code{back\_b} (ein \code{back\_b\_in}, aus \code{back\_b\_out}).

\begin{referenz}{@{}>{\raggedright\arraybackslash}p{0.2\linewidth}>{\raggedright\arraybackslash}p{0.11\linewidth}>{\raggedright\arraybackslash}p{0.13\linewidth}>{\raggedright\arraybackslash}p{0.15\linewidth}>{\raggedright\arraybackslash}X@{}}
\toprule
\textbf{Parameter} & \textbf{Typ} & \textbf{Standard} & \textbf{Einheit / Bereich} & \textbf{Bedeutung} \\
\midrule
\endhead
\param{point} & Vektor & \code{0.0, 0.0, 0.0} & mm & Bezugspunkt der Ebene bzw.\ des Elements. \\
\param{normal} & Vektor & \code{0.0, 0.0, -1.0} & -- & Flächennormale. \\
\param{up} & Vektor & \code{0.0, 1.0, 0.0} & -- & Richtung der lokalen v-Achse. \\
\param{drive} & Zahl & \code{0.0} & mm & Verschiebung längs der Normalen. \\
\param{r} & Zahl & $1/\sqrt{2}$ & -- & Amplitudenreflexionsgrad. \\
\param{t} & Zahl & $1/\sqrt{2}$ & -- & Amplitudentransmissionsgrad. \\
\bottomrule
\end{referenz}

\subsection{\block{fg.grating}}\label{sec:block-fg-grating}\index{fg.grating@\texttt{fg.grating} (Block)}

\textbf{Kategorie:} \code{splitter} \quad \textbf{Anbieter:} \plugin{FourierGrating} (gerechnet in seinem Prozess)

\textit{Periodic sample (FourierGrating table): diffraction orders as channels}

Eine periodische Probe (Gitter): Jede Beugungsordnung wird ein eigener Kanal. Die Wirkung kommt
entweder aus einer gespeicherten Jones-Tabelle (rechnet ohne Löser) oder aus einer
FourierGrating-Konfiguration, aus der der Block seine Tabelle um die einfallenden Strahlen herum
selbst baut (R) bzw.\ jede ebene Welle eines Felds rechnet (H). \gui{Check workflow} speichert
Tabelle und Konfiguration im Block und leert die Pfade. Der frühere Typname
\block{fusion.grating} wird beim Öffnen auf \block{fg.grating} umgesetzt. Strahlen, deren Ordnung
verschwindet oder die außerhalb der Tabelle bzw.\ nahe einer Rayleigh-Anomalie liegen, werden mit
einer Log-Zeile verworfen.

\textbf{Licht:} \code{in} (ein, RHS, vorn); \code{out} (aus, RHS, vorn, \enquote{orders});
\code{back\_out} (aus, RHS, hinten, \enquote{transmitted orders}).

\textbf{Ports:} \code{front} (ein \code{in}, aus \code{out}); \code{back} (aus \code{back\_out}).
Streutabelle: \code{front} $\to$ \code{front} und \code{front} $\to$ \code{back}.

\code{out} gibt die Ordnungen der mit \param{side} gewählten Seite weiter (wie bisher, ein
Port). \code{back\_out} gibt immer die durchgelassenen Ordnungen: So liegt ein
Transmissionsgitter im Modell als Schritt zwischen zwei Nachbarn, etwa zwei Gitter hintereinander
(Licht in \code{in}, weiter aus \code{back\_out}). Eine Tabelle nur reflektierter Ordnungen am
hinteren Port meldet der Block (RTT-Fusion ab 0.5.1).

\begin{referenz}{@{}>{\raggedright\arraybackslash}p{0.2\linewidth}>{\raggedright\arraybackslash}p{0.11\linewidth}>{\raggedright\arraybackslash}p{0.13\linewidth}>{\raggedright\arraybackslash}p{0.15\linewidth}>{\raggedright\arraybackslash}X@{}}
\toprule
\textbf{Parameter} & \textbf{Typ} & \textbf{Standard} & \textbf{Einheit / Bereich} & \textbf{Bedeutung} \\
\midrule
\endhead
\param{table} & Datei & leer & -- & Jones-Tabelle (\code{raytatouille-result}-Datei); \gui{Check workflow} speichert sie im Block (\param{table\_data}) und leert den Pfad. \\
\param{fouriergrating} & Datei & leer & -- & FourierGrating-Konfiguration: Der Block baut seine Tabelle um die Strahlen (R) oder rechnet jede ebene Welle eines Felds (H); wird beim Prüfen im Block gespeichert (\param{config}). \\
\param{config} & Daten & -- & -- & Die im Block gehaltene Konfiguration (hat Vorrang vor der Datei). \\
\param{table\_\allowbreak{}data} & Daten & -- & -- & Die im Block gehaltene Jones-Tabelle (hat Vorrang vor der Datei). \\
\param{side} & Auswahl & \code{reflected} & \code{reflected}, \code{transmitted} & Ordnungen an \code{out} (mit einer Konfiguration): reflektiert oder durchgelassen; \code{back\_out} gibt immer die durchgelassenen. \\
\param{orders} & Zahlenliste & leer & -- & Weitergegebene Ordnungen; leer: alle der Tabelle. \\
\param{n\_\allowbreak{}in} & Zahl & \code{1.0} & -- & Brechzahl des Einfallsmediums (R, selbst gebaute Tabelle). \\
\param{sigma\_\allowbreak{}step} \newline\small\gui{$\sigma$ step of the table} & Zahl & \code{0.01} & $\geq$ $10^{-5}$ & Abstand der Stützstellen der Jones-Tabelle in den Richtungskosinus $\sigma = n \sin\theta$ (kubisch interpoliert); kleiner ist genauer und braucht mehr FourierGrating-Rechnungen. \\
\param{sigma\_\allowbreak{}margin} \newline\small\gui{$\sigma$ margin of the table} & Zahl & \code{0.01} & $\geq$ 0.0 & Wie weit die Tabelle über den $\sigma$-Bereich der einfallenden Strahlen hinausreicht. \\
\param{point} & Vektor & \code{0.0, 0.0, 0.0} & mm & Bezugspunkt des Gitters (Phase bezogen auf $x = 0$). \\
\param{normal} & Vektor & \code{0.0, 0.0, 1.0} & -- & Flächennormale des Gitters. \\
\param{grating\_\allowbreak{}vector} & Vektor & \code{1.0, 0.0, 0.0} & -- & Gittervektor: lokale x-Achse, senkrecht zu den Furchen. \\
\param{field\_\allowbreak{}transfer} & Auswahl & \code{exact} & \code{exact}, \code{interpolated} & \code{exact}: FourierGrating je ebener Welle; \code{interpolated}: aus einem geprüften Teilraster des Spektrums ($10^{-8}$), viel schneller für breite Bündel. \\
\param{field\_\allowbreak{}threshold} & Zahl & $10^{-10}$ & -- & Ebene Wellen unter diesem Anteil der größten Amplitude werden weggelassen (H). \\
\param{field\_\allowbreak{}grid\_\allowbreak{}max} & Ganzzahl & \code{1024} & -- & Höchstens so viele Gitterpunkte je Achse für ein Feld (H). \\
\bottomrule
\end{referenz}

\subsection{\block{hpfem.grating}}\label{sec:block-hpfem-grating}\index{hpfem.grating@\texttt{hpfem.grating} (Block)}

\textbf{Kategorie:} \code{splitter} \quad \textbf{Anbieter:} \plugin{hp-FEM} (gerechnet in seinem Prozess; braucht \code{rtt-fusion} ab 0.5)

\textit{Periodic sample computed by hp-FEM (finite elements): diffraction orders as channels}

Eine periodische Probe wie \block{fg.grating}, gerechnet mit hp-FEM: Jede Beugungsordnung wird ein
eigener Kanal. Das Gitter ist ein hp-FEM-Job (Zelle, Netz, Materialien, Stapel, Löser) ohne
Einfall, Sweep und Karten; den Einfall setzen die Strahlen. Um die Strahlen einer Wellenlänge
baut der Block eine Jones-Tabelle: für eine Richtung genau dort, sonst Stützstellen im Abstand
\param{sigma\_step}, höchstens \param{max\_nodes}, je Stützstelle zwei hp-FEM-Lösungen. Der Block
rechnet in Strahlen (nativ R); Felder und Spektren werden vorher zu Strahlen. \gui{Check workflow}
speichert den Job im Block (\param{job\_data}) und leert den Pfad; eine Netzdatei bleibt außerhalb
und wird gemeldet. Bedienung: \autoref{sec:hp-block}.

\textbf{Licht:} \code{in} (ein, RHS, vorn); \code{out} (aus, RHS, vorn, \enquote{orders});
\code{back\_out} (aus, RHS, hinten, \enquote{transmitted orders}).

\textbf{Ports:} \code{front} (ein \code{in}, aus \code{out}); \code{back} (aus \code{back\_out}).
Streutabelle: \code{front} $\to$ \code{front} und \code{front} $\to$ \code{back}.

\code{out} gibt die Ordnungen der mit \param{side} gewählten Seite weiter (wie bisher, ein
Port). \code{back\_out} gibt immer die durchgelassenen Ordnungen: So liegt ein
Transmissionsgitter im Modell als Schritt zwischen zwei Nachbarn, etwa zwei Gitter hintereinander
(Licht in \code{in}, weiter aus \code{back\_out}). Eine Tabelle nur reflektierter Ordnungen am
hinteren Port meldet der Block (RTT-Fusion ab 0.5.1).

\begin{referenz}{@{}>{\raggedright\arraybackslash}p{0.2\linewidth}>{\raggedright\arraybackslash}p{0.11\linewidth}>{\raggedright\arraybackslash}p{0.13\linewidth}>{\raggedright\arraybackslash}p{0.15\linewidth}>{\raggedright\arraybackslash}X@{}}
\toprule
\textbf{Parameter} & \textbf{Typ} & \textbf{Standard} & \textbf{Einheit / Bereich} & \textbf{Bedeutung} \\
\midrule
\endhead
\param{job} & Datei & leer & -- & hp-FEM-Job (JSON, Problem \code{grating}); \gui{Check workflow} speichert ihn im Block (\param{job\_data}) und leert den Pfad. \\
\param{job\_\allowbreak{}data} & Daten & -- & -- & Der im Block gehaltene Job (hat Vorrang vor der Datei); \gui{Store job in block} setzt ihn. \\
\param{side} & Auswahl & \code{reflected} & \code{reflected}, \code{transmitted} & Ordnungen an \code{out}: reflektiert oder durchgelassen; \code{back\_out} gibt immer die durchgelassenen. \\
\param{orders} & Zahlenliste & leer & -- & Weitergegebene Ordnungen; leer: $-1$, 0, 1. \\
\param{sigma\_\allowbreak{}step} \newline\small\gui{$\sigma$ step of the table} & Zahl & \code{0.01} & $\geq$ $10^{-5}$ & Abstand der Stützstellen in den Richtungskosinus $\sigma = n \sin\theta$ (kubisch interpoliert); jede Stützstelle sind zwei hp-FEM-Lösungen. \\
\param{sigma\_\allowbreak{}margin} \newline\small\gui{$\sigma$ margin of the table} & Zahl & \code{0.01} & $\geq$ 0.0 & Wie weit die Tabelle über den $\sigma$-Bereich der Strahlen hinausreicht. \\
\param{max\_\allowbreak{}nodes} \newline\small\gui{Most table nodes} & Ganzzahl & \code{400} & $\geq$ 16 & Ein breiteres Bündel bricht mit einer Meldung ab, statt so viele Lösungen zu starten. \\
\param{point} & Vektor & \code{0.0, 0.0, 0.0} & mm & Oberkante des Gitters; dort liegt $x = 0$ des Jobs. \\
\param{normal} & Vektor & \code{0.0, 0.0, 1.0} & -- & Flächennormale (das $y$ des Jobs). \\
\param{grating\_\allowbreak{}vector} & Vektor & \code{1.0, 0.0, 0.0} & -- & Lokale x-Achse (das $x$ des Jobs), senkrecht zu den Linien. \\
\bottomrule
\end{referenz}

\section{Detektoren}\label{sec:bloecke-detektoren}

\subsection{\block{rtt.detector}}\label{sec:block-rtt-detector}\index{rtt.detector@\texttt{rtt.detector} (Block)}

\textbf{Kategorie:} \code{detector} \quad \textbf{Anbieter:} \plugin{Raytracer} (gerechnet in seinem Prozess; ohne Raytracer rechnet die Workflow-Maschine selbst, mit \code{raytatouille})

\textit{Detector / camera: amplitudes within, intensities across coherence groups}

Detektor bzw.\ Kamera eines Raytatouille-Systems: Amplituden addieren sich innerhalb einer
Kohärenzgruppe, Intensitäten über die Gruppen. Der Ausgang \code{image} liefert das Bild an
einen \block{display.image}. Ist im Workflow-Panel ein \gui{Image window} größer als 0
eingestellt, ersetzt es \param{half\_size}.

\textbf{Licht:} \code{in} (ein, RHGSP).\\
\textbf{Daten:} \code{image} (aus, \code{result}).

\textbf{Ports:} \code{in} (ein \code{in}).

\begin{referenz}{@{}>{\raggedright\arraybackslash}p{0.2\linewidth}>{\raggedright\arraybackslash}p{0.11\linewidth}>{\raggedright\arraybackslash}p{0.13\linewidth}>{\raggedright\arraybackslash}p{0.15\linewidth}>{\raggedright\arraybackslash}X@{}}
\toprule
\textbf{Parameter} & \textbf{Typ} & \textbf{Standard} & \textbf{Einheit / Bereich} & \textbf{Bedeutung} \\
\midrule
\endhead
\param{surface} & Text & leer & -- & Kennung der Fläche im Raytatouille-System. \\
\param{pixels} \newline\small\gui{pixels (long side)} & Ganzzahl & \code{512} & 16 … 8192 & Pixelzahl an der langen Seite. \\
\param{half\_\allowbreak{}size} & Zahlenliste & leer & mm & Bildfenster $[h_u, h_v]$ (halbe Breiten); leer: die Apertur des Detektors. \\
\bottomrule
\end{referenz}

\subsection{\block{rtt.part.detector}}\label{sec:block-rtt-part-detector}\index{rtt.part.detector@\texttt{rtt.part.detector} (Block)}

\textbf{Kategorie:} \code{detector} \quad \textbf{Anbieter:} \plugin{Raytracer} (gerechnet in seinem Prozess; ohne Raytracer rechnet die Workflow-Maschine selbst, mit \code{raytatouille})

\textit{Raytatouille detector (a new element)}

Ein neuer Raytatouille-Detektor aus der Bibliothek. Vorlage: rechteckige Fläche
$10 \times 10$~mm. \gui{New Raytatouille system} legt einen solchen Detektor bei $z = 50$~mm an.

\textbf{Licht:} \code{in\_in} (ein, RHS, vorn).\\
\textbf{Daten:} \code{image} (aus, \code{result}).

\textbf{Ports:} \code{in} (ein \code{in\_in}).

\begin{referenz}{@{}>{\raggedright\arraybackslash}p{0.2\linewidth}>{\raggedright\arraybackslash}p{0.11\linewidth}>{\raggedright\arraybackslash}p{0.13\linewidth}>{\raggedright\arraybackslash}p{0.15\linewidth}>{\raggedright\arraybackslash}X@{}}
\toprule
\textbf{Parameter} & \textbf{Typ} & \textbf{Standard} & \textbf{Einheit / Bereich} & \textbf{Bedeutung} \\
\midrule
\endhead
\param{rtt\_\allowbreak{}new} & Daten & -- & -- & Das Raytatouille-Element (Lage, Flächen, Glas); im Dialog als eigene Felder. \\
\bottomrule
\end{referenz}

\subsection{\block{optics.detector}}\label{sec:block-optics-detector}\index{optics.detector@\texttt{optics.detector} (Block)}

\textbf{Kategorie:} \code{detector} \quad \textbf{Anbieter:} \plugin{Workflow} (RTT-Fusion, eingebaut)

\textit{Plane detector: coherent image (amplitudes per coherence group)}

Ein ebener Detektor (analytisch): kohärentes Bild, Amplituden je Kohärenzgruppe.

\textbf{Licht:} \code{in} (ein, RHS).\\
\textbf{Daten:} \code{image} (aus, \code{result}).

\textbf{Ports:} \code{in} (ein \code{in}).

\begin{referenz}{@{}>{\raggedright\arraybackslash}p{0.2\linewidth}>{\raggedright\arraybackslash}p{0.11\linewidth}>{\raggedright\arraybackslash}p{0.13\linewidth}>{\raggedright\arraybackslash}p{0.15\linewidth}>{\raggedright\arraybackslash}X@{}}
\toprule
\textbf{Parameter} & \textbf{Typ} & \textbf{Standard} & \textbf{Einheit / Bereich} & \textbf{Bedeutung} \\
\midrule
\endhead
\param{point} & Vektor & \code{0.0, 0.0, 0.0} & mm & Bezugspunkt der Ebene bzw.\ des Elements. \\
\param{normal} & Vektor & \code{0.0, 0.0, -1.0} & -- & Flächennormale. \\
\param{up} & Vektor & \code{0.0, 1.0, 0.0} & -- & Richtung der lokalen v-Achse. \\
\param{drive} & Zahl & \code{0.0} & mm & Verschiebung längs der Normalen. \\
\param{half\_\allowbreak{}size} & Zahlenliste & \code{1.0, 1.0} & mm & Halbe Breiten des Bildfensters $[h_u, h_v]$. \\
\param{pixels} & Zahlenliste & \code{128, 128} & -- & Pixel in u und v. \\
\bottomrule
\end{referenz}

\section{Umwandlung, Auswertung und Daten}\label{sec:bloecke-auswertung}


\begin{hinweis}
Für \block{analysis.camera\_image} und \block{biooptomd.layer} hat die Workflow-Maschine in der
vorliegenden Version keinen Rechner, für \block{fusion.convert} erst ab RTT-Fusion 0.2.1.
\gui{Add block} bietet solche Typen nicht an; steht ein solcher Block schon im Workflow, warnt
\gui{Check workflow} mit \code{check.no\_evaluator}, und ein Lauf, der ihn braucht, bricht mit
\code{exec.no\_evaluator} ab. Umwandlungen plant der Planer ohnehin selbst; für Kamerabilder
verwenden Sie den Detektor eines Systems mit einem \block{display.image}.
\end{hinweis}


\subsection{\block{fusion.convert}}\label{sec:block-fusion-convert}\index{fusion.convert@\texttt{fusion.convert} (Block)}

\textbf{Kategorie:} \code{conversion} \quad \textbf{Anbieter:} \plugin{Workflow} (RTT-Fusion, eingebaut)

\textit{Conversion between representations (R $\to$ H, H $\leftrightarrow$ G, G $\leftrightarrow$ S, …)}

Eine ausdrückliche Umwandlung zwischen Darstellungen (R $\to$ H, H $\leftrightarrow$ G,
G $\leftrightarrow$ S …). Normalerweise plant der Planer Umwandlungen selbst; sichtbar wird eine
Umwandlung nur, wenn Sie sie festlegen.

Gerechnet wird ab RTT-Fusion 0.2.1 (mit \code{fieldtracing}): R $\to$ H (Strahlen zu einem Feld
auf der Portebene) und H $\to$ R, mit Spektren (RTT-Fusion R30) auch \param{to}~=~\code{S}:
Strahlen über ein Feld zum Spektrum ebener Wellen (R $\to$ H $\to$ S), ein Feld über die PFT oder
eine FFT (H $\to$ S), und zurück. Licht, das schon in der Zieldarstellung ankommt, geht
unverändert durch. Ein Spektrum hinter dem Block bleibt auf einer Verbindung auf \code{auto} ein
Spektrum, wenn der nächste Block S nativ nimmt (z.\,B.\ eine Kamera); die Verbindung zeigt nach
dem Lauf dann \textbf{R$\to$S} (Abschnitt~\ref{sec:workflow-verwendet}). Die Ziele G und P kommen
in einer späteren Stufe; ein Lauf meldet dann \code{exec.convert\_unsupported}, ebenso wenn
\code{fieldtracing} fehlt. Standard ist \code{H}.

\textbf{Licht:} \code{in} (ein, RHGSTP); \code{out} (aus, RHGSTP).

\textbf{Ports:} \code{in} (ein \code{in}); \code{out} (aus \code{out}).

\begin{referenz}{@{}>{\raggedright\arraybackslash}p{0.2\linewidth}>{\raggedright\arraybackslash}p{0.11\linewidth}>{\raggedright\arraybackslash}p{0.13\linewidth}>{\raggedright\arraybackslash}p{0.15\linewidth}>{\raggedright\arraybackslash}X@{}}
\toprule
\textbf{Parameter} & \textbf{Typ} & \textbf{Standard} & \textbf{Einheit / Bereich} & \textbf{Bedeutung} \\
\midrule
\endhead
\param{to} & Auswahl & \code{H} & \code{R}, \code{H}, \code{G}, \code{S}, \code{P} & Ziel-Darstellung. \\
\bottomrule
\end{referenz}

\subsection{\block{analysis.camera\_image}}\label{sec:block-analysis-camera-image}\index{analysis.cameraimage@\texttt{analysis.camera\_image} (Block)}

\textbf{Kategorie:} \code{analysis} \quad \textbf{Anbieter:} \plugin{Workflow} (RTT-Fusion, eingebaut)

\textit{Irradiance image of a detector with many rays}

Bestrahlungsstärkebild eines Detektors mit vielen Strahlen.

\textbf{Licht:} \code{in} (ein, RHGSP).\\
\textbf{Daten:} \code{image} (aus, \code{result}).

\textbf{Ports:} \code{in} (ein \code{in}).

\begin{referenz}{@{}>{\raggedright\arraybackslash}p{0.2\linewidth}>{\raggedright\arraybackslash}p{0.11\linewidth}>{\raggedright\arraybackslash}p{0.13\linewidth}>{\raggedright\arraybackslash}p{0.15\linewidth}>{\raggedright\arraybackslash}X@{}}
\toprule
\textbf{Parameter} & \textbf{Typ} & \textbf{Standard} & \textbf{Einheit / Bereich} & \textbf{Bedeutung} \\
\midrule
\endhead
\param{rays} \newline\small\gui{rays per field and wavelength} & Ganzzahl & \code{200000} & $\geq$ 1000 & Strahlen je Feld und Wellenlänge. \\
\param{pixels} & Ganzzahl & \code{512} & 16 … 8192 & Pixelzahl. \\
\bottomrule
\end{referenz}

\subsection{\block{biooptomd.layer}}\label{sec:block-biooptomd-layer}\index{biooptomd.layer@\texttt{biooptomd.layer} (Block)}

\textbf{Kategorie:} \code{data} \quad \textbf{Anbieter:} \plugin{Workflow} (RTT-Fusion, eingebaut)

\textit{Refractive-index profile n(z) and dipole ensembles from BioOptoMD}

Datenquelle: Brechzahlprofil $n(z)$ (Datentyp \code{N}) und Dipol-Ensembles (Datentyp \code{Q})
aus einem Ausgabeordner von BioOptoMD.

\textbf{Daten:} \code{profile} (aus, \code{N}); \code{dipoles} (aus, \code{Q}).

\begin{referenz}{@{}>{\raggedright\arraybackslash}p{0.2\linewidth}>{\raggedright\arraybackslash}p{0.11\linewidth}>{\raggedright\arraybackslash}p{0.13\linewidth}>{\raggedright\arraybackslash}p{0.15\linewidth}>{\raggedright\arraybackslash}X@{}}
\toprule
\textbf{Parameter} & \textbf{Typ} & \textbf{Standard} & \textbf{Einheit / Bereich} & \textbf{Bedeutung} \\
\midrule
\endhead
\param{directory} \newline\small\gui{BioOptoMD output} & Datei & leer & -- & Ausgabeordner von BioOptoMD. \\
\bottomrule
\end{referenz}

\subsection{Materialblöcke des Path Tracers}\label{sec:block-pathtracer-material}\index{pathtracer.material@\texttt{pathtracer.material} (Block)}

\block{pathtracer.bsdf\_file}, \block{pathtracer.bsdf\_analytic}, \block{pathtracer.bsdf\_mix} und
\block{pathtracer.material} haben nur Datenanschlüsse vom Typ \code{bsdf} und kein Licht; ebenso
\block{pathtracer.nk}, \block{pathtracer.topography} und \block{pathtracer.rough\_surface} mit den
Typen \code{nk} und \code{topography}. Sie
beschreiben Materialien, die der Path Tracer beim Rechnen selbst aus dem Workflow liest; Parameter
und Bedeutung in \autoref{sec:pt-materialbloecke} und \autoref{sec:ref-pt-bloecke}.

\section{Anzeigeblöcke}\label{sec:bloecke-anzeige}


Anzeigeblöcke bestimmen, was \gui{Run workflow} rechnet (Abschnitt~\ref{sec:workflow-was}). Ein
Doppelklick öffnet ihr Fenster (Abschnitt~\ref{sec:workflow-anzeigen}). Ihr Ausgang \code{view}
ist das, was das Fenster zeigt.


\subsection{\block{display.image}}\label{sec:block-display-image}\index{display.image@\texttt{display.image} (Block)}

\textbf{Kategorie:} \code{analysis} \quad \textbf{Anbieter:} \plugin{Workflow} (RTT-Fusion, eingebaut)

\textit{Display: an image result (camera, detector) in its own window}

Zeigt ein Bild (Kamera, Detektor) in einem eigenen Fenster. Eingang: der Ausgang \code{image}
eines Detektors.

\textbf{Daten:} \code{in} (ein, \code{result}, \enquote{image}); \code{view} (aus, \code{result}).

\begin{referenz}{@{}>{\raggedright\arraybackslash}p{0.2\linewidth}>{\raggedright\arraybackslash}p{0.11\linewidth}>{\raggedright\arraybackslash}p{0.13\linewidth}>{\raggedright\arraybackslash}p{0.15\linewidth}>{\raggedright\arraybackslash}X@{}}
\toprule
\textbf{Parameter} & \textbf{Typ} & \textbf{Standard} & \textbf{Einheit / Bereich} & \textbf{Bedeutung} \\
\midrule
\endhead
\param{title} & Text & leer & -- & Fenstertitel; leer: der Name des Blocks. \\
\bottomrule
\end{referenz}

\subsection{\block{display.plot}}\label{sec:block-display-plot}\index{display.plot@\texttt{display.plot} (Block)}

\textbf{Kategorie:} \code{analysis} \quad \textbf{Anbieter:} \plugin{Workflow} (RTT-Fusion, eingebaut)

\textit{Display: curves of a sweep / optimisation, or profiles of an image}

Zeigt Kurven: an einem Meta-Block die Messgröße über dem variierten Parameter, an einem Bild seine
mittlere Zeile und Spalte, an Zahlen ihre Werte über dem Index.

\textbf{Daten:} \code{in} (ein, \code{result}, \enquote{result}); \code{view} (aus, \code{result}).

\begin{referenz}{@{}>{\raggedright\arraybackslash}p{0.2\linewidth}>{\raggedright\arraybackslash}p{0.11\linewidth}>{\raggedright\arraybackslash}p{0.13\linewidth}>{\raggedright\arraybackslash}p{0.15\linewidth}>{\raggedright\arraybackslash}X@{}}
\toprule
\textbf{Parameter} & \textbf{Typ} & \textbf{Standard} & \textbf{Einheit / Bereich} & \textbf{Bedeutung} \\
\midrule
\endhead
\param{title} & Text & leer & -- & Fenstertitel; leer: der Name des Blocks. \\
\bottomrule
\end{referenz}

\subsection{\block{display.beam}}\label{sec:block-display-beam}\index{display.beam@\texttt{display.beam} (Block)}

\textbf{Kategorie:} \code{analysis} \quad \textbf{Anbieter:} \plugin{Workflow} (RTT-Fusion, eingebaut)

\textit{Display: the light at a block output (probe), intensity on its port plane}

Eine Sonde für das Licht an einem beliebigen Ausgang: Intensität auf der Port-Ebene (durch den
Schwerpunkt des Lichts, senkrecht zu seiner mittleren Richtung). Strahlen werden als kohärentes
Strahlbild, Felder als Feldbild gezeigt; Amplituden addieren sich innerhalb einer
Kohärenzgruppe, Intensitäten über die Gruppen. Verbunden wird sie, indem Sie einen Port auf den
Block ziehen (Beobachter-Link).

\textbf{Licht:} \code{in} (ein, RH).\\
\textbf{Daten:} \code{view} (aus, \code{result}).

\textbf{Ports:} \code{in} (ein \code{in}).

\begin{referenz}{@{}>{\raggedright\arraybackslash}p{0.2\linewidth}>{\raggedright\arraybackslash}p{0.11\linewidth}>{\raggedright\arraybackslash}p{0.13\linewidth}>{\raggedright\arraybackslash}p{0.15\linewidth}>{\raggedright\arraybackslash}X@{}}
\toprule
\textbf{Parameter} & \textbf{Typ} & \textbf{Standard} & \textbf{Einheit / Bereich} & \textbf{Bedeutung} \\
\midrule
\endhead
\param{title} & Text & leer & -- & Fenstertitel; leer: der Name des Blocks. \\
\param{pixels} & Ganzzahl & \code{128} & 16 … 2048 & Pixelzahl des Sondenbilds. \\
\param{margin} & Zahl & \code{0.1} & $\geq$ 0.0 & Rand des Fensters um das Licht (relativ). \\
\bottomrule
\end{referenz}

\section{Meta-Blöcke}\label{sec:bloecke-meta}


Meta-Blöcke wiederholen den Teil des Workflows vor einem Ausgang mit geänderten Parametern
(Abschnitt~\ref{sec:workflow-meta}). Alle drei haben dieselben Anschlüsse: \code{observe}
(der beobachtete Ausgang, alternativ zum Parameter \param{collect}), \code{drive} (Antrieb in den
\code{drive}- oder \code{pose}-Eingang eines Blocks, alternativ zum Parameter \param{parameter})
und \code{results} (die gesammelten Ergebnisse, z.\,B.\ für \block{display.plot}). Messgrößen
eines Bildes: \code{peak} (Maximum), \code{centre} (Mittel der $3 \times 3$ mittleren Pixel),
\code{power} (Summe); eine Zahl ist ihre eigene Messgröße.


\subsection{\block{meta.sweep}}\label{sec:block-meta-sweep}\index{meta.sweep@\texttt{meta.sweep} (Block)}

\textbf{Kategorie:} \code{meta} \quad \textbf{Anbieter:} \plugin{Workflow} (RTT-Fusion, eingebaut)

\textit{Sweep: a parameter over values, the output collected per value}

Ein Parameter über eine Liste von Werten; die Ausgabe wird je Wert gesammelt.

\textbf{Daten:} \code{results} (aus, \code{result}); \code{drive} (aus, \code{drive}); \code{observe} (ein, \code{result}).

\begin{referenz}{@{}>{\raggedright\arraybackslash}p{0.2\linewidth}>{\raggedright\arraybackslash}p{0.11\linewidth}>{\raggedright\arraybackslash}p{0.13\linewidth}>{\raggedright\arraybackslash}p{0.15\linewidth}>{\raggedright\arraybackslash}X@{}}
\toprule
\textbf{Parameter} & \textbf{Typ} & \textbf{Standard} & \textbf{Einheit / Bereich} & \textbf{Bedeutung} \\
\midrule
\endhead
\param{parameter} & Text & leer & -- & Variierter Parameter als \code{Block.Parameter} oder \code{Block.Parameter[i]}, z.\,B.\ \code{mirror.point[0]}. \\
\param{collect} & Text & leer & -- & Gesammelter Ausgang als \code{Block.Anschluss}, z.\,B.\ \code{camera.image}. \\
\param{measure} & Auswahl & \code{peak} & \code{peak}, \code{centre}, \code{power} & Messgröße jedes Ergebnisses. \\
\param{component} & Auswahl & \code{dz} & \code{dx}, \code{dy}, \code{dz}, \code{rx}, \code{ry}, \code{rz} & Komponente der Lageänderung, wenn \code{drive} in einen \code{pose}-Eingang geht (mm bzw.\ rad). \\
\param{values} & Zahlenliste & leer & -- & Werte; leer: aus \param{start}, \param{stop}, \param{count}. \\
\param{start} & Zahl & \code{0.0} & -- & Erster Wert. \\
\param{stop} & Zahl & \code{1.0} & -- & Letzter Wert. \\
\param{count} & Ganzzahl & \code{11} & $\geq$ 1 & Anzahl der Werte. \\
\bottomrule
\end{referenz}

\subsection{\block{meta.replicates}}\label{sec:block-meta-replicates}\index{meta.replicates@\texttt{meta.replicates} (Block)}

\textbf{Kategorie:} \code{meta} \quad \textbf{Anbieter:} \plugin{Workflow} (RTT-Fusion, eingebaut)

\textit{Replicates: repeated with replicate values, images averaged}

Wiederholungen mit Wiederholungswerten (z.\,B.\ Zufalls-Startwerte, Realisierungen); die
gesammelten Bilder werden gemittelt.

\textbf{Daten:} \code{results} (aus, \code{result}); \code{drive} (aus, \code{drive}); \code{observe} (ein, \code{result}).

\begin{referenz}{@{}>{\raggedright\arraybackslash}p{0.2\linewidth}>{\raggedright\arraybackslash}p{0.11\linewidth}>{\raggedright\arraybackslash}p{0.13\linewidth}>{\raggedright\arraybackslash}p{0.15\linewidth}>{\raggedright\arraybackslash}X@{}}
\toprule
\textbf{Parameter} & \textbf{Typ} & \textbf{Standard} & \textbf{Einheit / Bereich} & \textbf{Bedeutung} \\
\midrule
\endhead
\param{parameter} & Text & leer & -- & Variierter Parameter als \code{Block.Parameter} oder \code{Block.Parameter[i]}, z.\,B.\ \code{mirror.point[0]}. \\
\param{collect} & Text & leer & -- & Gesammelter Ausgang als \code{Block.Anschluss}, z.\,B.\ \code{camera.image}. \\
\param{measure} & Auswahl & \code{peak} & \code{peak}, \code{centre}, \code{power} & Messgröße jedes Ergebnisses. \\
\param{component} & Auswahl & \code{dz} & \code{dx}, \code{dy}, \code{dz}, \code{rx}, \code{ry}, \code{rz} & Komponente der Lageänderung, wenn \code{drive} in einen \code{pose}-Eingang geht (mm bzw.\ rad). \\
\param{values} & Zahlenliste & leer & -- & Wiederholungswerte; leer: $0 \ldots \text{count} - 1$. \\
\param{count} & Ganzzahl & \code{8} & $\geq$ 1 & Anzahl der Wiederholungen (wenn \param{values} leer ist). \\
\bottomrule
\end{referenz}

\subsection{\block{meta.optimise}}\label{sec:block-meta-optimise}\index{meta.optimise@\texttt{meta.optimise} (Block)}

\textbf{Kategorie:} \code{meta} \quad \textbf{Anbieter:} \plugin{Workflow} (RTT-Fusion, eingebaut)

\textit{Optimisation of one parameter for the extremum of a measure}

Ein Parameter innerhalb von Grenzen wird so gesucht, dass die Messgröße des gesammelten Ergebnisses
ihr Maximum oder Minimum erreicht: erst ein grobes Raster, dann eine eindimensionale Suche.

\textbf{Daten:} \code{results} (aus, \code{result}); \code{drive} (aus, \code{drive}); \code{observe} (ein, \code{result}).

\begin{referenz}{@{}>{\raggedright\arraybackslash}p{0.2\linewidth}>{\raggedright\arraybackslash}p{0.11\linewidth}>{\raggedright\arraybackslash}p{0.13\linewidth}>{\raggedright\arraybackslash}p{0.15\linewidth}>{\raggedright\arraybackslash}X@{}}
\toprule
\textbf{Parameter} & \textbf{Typ} & \textbf{Standard} & \textbf{Einheit / Bereich} & \textbf{Bedeutung} \\
\midrule
\endhead
\param{parameter} & Text & leer & -- & Variierter Parameter als \code{Block.Parameter} oder \code{Block.Parameter[i]}, z.\,B.\ \code{mirror.point[0]}. \\
\param{collect} & Text & leer & -- & Gesammelter Ausgang als \code{Block.Anschluss}, z.\,B.\ \code{camera.image}. \\
\param{measure} & Auswahl & \code{peak} & \code{peak}, \code{centre}, \code{power} & Messgröße jedes Ergebnisses. \\
\param{component} & Auswahl & \code{dz} & \code{dx}, \code{dy}, \code{dz}, \code{rx}, \code{ry}, \code{rz} & Komponente der Lageänderung, wenn \code{drive} in einen \code{pose}-Eingang geht (mm bzw.\ rad). \\
\param{bounds} & Zahlenliste & \code{0.0, 1.0} & -- & Untere und obere Grenze. \\
\param{goal} & Auswahl & \code{maximum} & \code{maximum}, \code{minimum} & Gesucht ist das Maximum oder Minimum. \\
\param{scan} & Ganzzahl & \code{9} & $\geq$ 3 & Punkte des groben Rasters vor der Suche. \\
\param{tolerance} & Zahl & \code{0.0} & -- & Abbruchgenauigkeit; 0: $10^{-6}$ der Grenzen. \\
\bottomrule
\end{referenz}

\subsection{\block{fusion.actuator}}\label{sec:block-fusion-actuator}\index{fusion.actuator@\texttt{fusion.actuator} (Block)}

\textbf{Kategorie:} \code{meta} \quad \textbf{Anbieter:} \plugin{Workflow} (RTT-Fusion, eingebaut)

\textit{Actuator: a stage that sets a system parameter and carries blocks along}

Ein Tisch des Aufbaus (Verschiebetisch, Drehhalter, Fokusantrieb) mit einem Wert
(Abschnitt~\ref{sec:workflow-aktuatoren}). Der Wert setzt einen freien Parameter des
Raytatouille-Systems und verschiebt die Blöcke, die der Tisch trägt. Der Block rechnet nichts
selbst; die Workflow-Maschine wendet ihn vor jedem Lauf an.

\textbf{Daten:} \code{drive} (ein, \code{drive}).

\begin{referenz}{@{}>{\raggedright\arraybackslash}p{0.2\linewidth}>{\raggedright\arraybackslash}p{0.11\linewidth}>{\raggedright\arraybackslash}p{0.13\linewidth}>{\raggedright\arraybackslash}p{0.15\linewidth}>{\raggedright\arraybackslash}X@{}}
\toprule
\textbf{Parameter} & \textbf{Typ} & \textbf{Standard} & \textbf{Einheit / Bereich} & \textbf{Bedeutung} \\
\midrule
\endhead
\param{value} & Zahl & \code{0.0} & -- & Stellung des Tischs. \\
\param{unit} & Auswahl & \code{mm} & \code{mm}, \code{deg}, \code{rad}, \code{µm} & Einheit des Werts; \code{deg}: Drehungen der getragenen Blöcke in Grad. \\
\param{parameter} & Text & leer & -- & Freier Parameter des Raytatouille-Systems, der den Wert erhält, z.\,B.\ \code{DX}; innerhalb \code{min}/\code{max} seiner Zeile. \\
\param{moves} & Text & leer & -- & Getragene Blöcke: \code{Block.Komponente} oder \code{Block.Komponente*Faktor}, durch Kommas getrennt (\code{dx}, \code{dy}, \code{dz} in mm; \code{rx}, \code{ry}, \code{rz}), z.\,B.\ \code{G1.dx, G2.dx}. \\
\param{zero} & Zahl & \code{0.0} & -- & Wert, bei dem die getragenen Blöcke dort stehen, wo das Modell sie hinstellt. \\
\param{drive} & Zahl & \code{0.0} & -- & Wird zum Wert addiert (vom \code{drive} eines Meta-Blocks gesetzt). \\
\bottomrule
\end{referenz}

\section{Strukturblöcke}\label{sec:bloecke-struktur}


Diese Typen legen Sie nicht über \gui{Add block} an; sie entstehen durch Befehle des Editors oder
beim Abrollen. Ihre Anschlüsse ergeben sich aus ihrem Inhalt.


\subsection{\block{group}}\label{sec:block-group}\index{group@\texttt{group} (Block)}

\textbf{Kategorie:} \code{domain} \quad \textbf{Anbieter:} \plugin{Workflow} (RTT-Fusion, eingebaut)

\textit{Group (R28): a sub-model with variants; in the socket form it keeps the group's pose and exported parameters for the blocks inside}

Eine Gruppe (Abschnitt~\ref{sec:workflow-gruppen}): ein Teilmodell mit Varianten, angelegt mit
\gui{Group selected blocks…} oder \gui{Section $\to$ group}. Nach außen hat sie die exportierten
Ports der aktiven Variante; sie trägt die Lage der Gruppe und die nach außen gegebenen Parameter.

\textbf{Parameter:} keine.

\subsection{\block{fusion.condensed}}\label{sec:block-fusion-condensed}\index{fusion.condensed@\texttt{fusion.condensed} (Block)}

\textbf{Kategorie:} \code{domain} \quad \textbf{Anbieter:} \plugin{Workflow} (RTT-Fusion, eingebaut)

\textit{Condensed group (R29): plane-wave tables of its inner branches}

Eine verdichtete Gruppe (Häkchen \gui{Condensed (plane-wave tables)}): Je innerem Ast wird einmal
eine Tabelle über einfallende ebene Wellen gerechnet und dann angewendet. Einstellungen beim
Einschalten: \code{sigma\_max} 0{,}2 und \code{count} 5.

\textbf{Parameter:} keine.

\subsection{\block{fusion.sum}}\label{sec:block-fusion-sum}\index{fusion.sum@\texttt{fusion.sum} (Block)}

\textbf{Kategorie:} \code{splitter} \quad \textbf{Anbieter:} \plugin{Workflow} (RTT-Fusion, eingebaut)

\textit{Junction of unrolled branches: the sum of the light arriving at one state}

Ein Knotenpunkt abgerollter Äste: addiert das Licht, das aus mehreren Ästen am selben Zustand
ankommt. Entsteht nur intern beim Abrollen.

\textbf{Parameter:} keine.

